%! The Idea of Elimination — Unit 2, Lesson 2.1 — Slides (Beamer)
\documentclass[aspectratio=169, 11pt]{beamer}
\usepackage{linalg-beamer}   % provides \burgundyheader{} and \sectionlabel[color]{}
\newcommand{\xx}{\mathbf{x}}
\newcommand{\bb}{\mathbf{b}}

\begin{document}

% TITLE SLIDE (CourseName is not defined in beamer, so the name is literal)
{
\setbeamercolor{background canvas}{bg=burgundy}
\begin{frame}[plain]
  \vspace{2.8cm}\hspace{0.55cm}%
  \begin{minipage}{0.9\textwidth}
    {\color{goldacc}\bfseries\small LESSON 2.1}\\[0.5cm]
    {\color{white}\bfseries\LARGE The Idea of Elimination}\\[0.3cm]
    {\color{white}\bfseries\large Pivots, Multipliers, and Back-Substitution}\\[1.0cm]
    {\color{blushmid}\small Linear Algebra~$\cdot$~Unit 2: Solving Linear Equations $A\xx=\bb$}
  \end{minipage}
\end{frame}
}

% HOOK
\begin{frame}[t, plain]
  \burgundyheader{When ``just subtract'' isn't enough}
  \sectionlabel{Hook}
  \begin{columns}[T]
    \begin{column}{0.48\textwidth}
      \textbf{Lesson 2.0}
      \[
        \begin{aligned}
          x + 2y &= 5 \\
          3x + 2y &= 11
        \end{aligned}
      \]
      Both have a $2y$ --- one subtraction removes $y$.
    \end{column}
    \begin{column}{0.48\textwidth}
      \textbf{Today}
      \[
        \begin{aligned}
          2x + 3y &= 13 \\
          4x + 7y &= 27
        \end{aligned}
      \]
      $x$-terms are $2$ and $4$ --- a plain subtract leaves \emph{both}.
    \end{column}
  \end{columns}
  \vspace{6pt}
  \begin{block}{The question}
    What can we do to the first equation \emph{first} so that subtracting wipes out $x$?
  \end{block}
\end{frame}

% PIVOT + MULTIPLIER
\begin{frame}[t, plain]
  \burgundyheader{Pivot and multiplier}
  \sectionlabel[goldacc]{The rule}
  \begin{itemize}
    \setlength{\itemsep}{5pt}
    \item The \textbf{pivot} is the first coefficient we build on: the $2$ in front of $x$.
    \item The \textbf{multiplier} says how many pivot rows to subtract:
  \end{itemize}
  \[
    \ell = \frac{\text{entry to eliminate}}{\text{pivot}} = \frac{4}{2} = 2.
  \]
  \begin{block}{Elimination step: row 2 $-\,\ell\times$ row 1}
    $(4x+7y) - 2(2x+3y) = 27 - 2(13) \;\Longrightarrow\; y = 1.$ \quad The $x$-term is gone \emph{by design}.
  \end{block}
\end{frame}

% TRIANGULAR + BACK-SUB
\begin{frame}[t, plain]
  \burgundyheader{Triangular, then back up}
  \sectionlabel{Finish the job}
  One step leaves an \textbf{upper-triangular} system --- zeros below the pivot:
  \[
    \begin{aligned}
      2x + 3y &= 13 \\
      y &= 1
    \end{aligned}
  \]
  \textbf{Back-substitute} upward: $\;2x + 3(1) = 13 \Rightarrow x = 5.$
  \begin{block}{Check by rebuilding $\bb$}
    $5\begin{bsmallmatrix} 2\\4 \end{bsmallmatrix} + 1\begin{bsmallmatrix} 3\\7 \end{bsmallmatrix} = \begin{bsmallmatrix} 13\\27 \end{bsmallmatrix}$ \checkmark \quad Solution $(x,y)=(5,1)$.
  \end{block}
\end{frame}

% STAIRCASE 3x3
\begin{frame}[t, plain]
  \burgundyheader{Bigger systems, same staircase}
  \sectionlabel[goldacc]{Scaling up}
  \[
    \begin{aligned}
      x + y + z &= 6 \\
      2x + 3y + z &= 11 \\
      x + 3y + 6z &= 25
    \end{aligned}
    \qquad\longrightarrow\qquad
    \begin{aligned}
      x + y + z &= 6 \\
      y - z &= -1 \\
      7z &= 21
    \end{aligned}
  \]
  \begin{itemize}
    \setlength{\itemsep}{4pt}
    \item Round 1: pivot in row 1 clears the whole first column ($\ell_{21}=2,\ \ell_{31}=1$).
    \item Round 2: new pivot in row 2 clears below it ($\ell_{32}=2$).
    \item Back-substitute: $z=3,\ y=2,\ x=1$.
  \end{itemize}
\end{frame}

% ZERO PIVOT
\begin{frame}[t, plain]
  \burgundyheader{When a pivot is zero}
  \sectionlabel{Road signs}
  \[
    \begin{aligned}
      0x + 2y &= 4 \\
      3x + y &= 5
    \end{aligned}
    \qquad\xrightarrow{\ \text{swap rows}\ }\qquad
    \begin{aligned}
      3x + y &= 5 \\
      2y &= 4
    \end{aligned}
  \]
  \begin{itemize}
    \setlength{\itemsep}{5pt}
    \item Pivot slot is $0$? \textbf{Exchange rows} to bring a nonzero up, then continue.
    \item No row to swap in $\Rightarrow$ breakdown: $0=(\text{nonzero})$ is \textbf{no solution};
          $0=0$ is \textbf{infinitely many}.
  \end{itemize}
\end{frame}

% PREVIEW
\begin{frame}[t, plain]
  \burgundyheader{Where this is heading}
  \sectionlabel{Connections}
  \textbf{Today:} pivot $\to$ multiplier $\to$ triangular form $\to$ back-substitute.\\[8pt]
  \begin{itemize}
    \setlength{\itemsep}{4pt}
    \item \textbf{2.2} Elimination Matrices and Inverse Matrices --- each step is a matrix.
    \item \textbf{2.3} Matrix Computations and $A=LU$ --- the multipliers become $L$.
    \item \textbf{2.4} Permutations and Transposes --- row exchanges, done right.
  \end{itemize}
\end{frame}

\end{document}
